% Context Deficit Maps — Workshop Paper Draft
% Compile: pdflatex main && bibtex main && pdflatex main && pdflatex main
%
% NOTE: Swap \documentclass to icml2026 or neurips_2026 for submission.
% This uses a clean article class for drafting.

\documentclass[11pt]{article}
\usepackage[margin=1in]{geometry}
\usepackage{amsmath,amssymb,amsthm}
\usepackage{booktabs}
\usepackage{graphicx}
\usepackage{hyperref}
\usepackage{xcolor}
\usepackage{natbib}
\usepackage{microtype}
\usepackage{caption}
\usepackage{subcaption}

% Clean paragraph spacing (no indent, vertical gap between paragraphs)
\usepackage{parskip}
\setlength{\parskip}{0.6\baselineskip plus 2pt minus 1pt}

% Draft helpers
\newcommand{\todo}[1]{\textcolor{red}{\textbf{[TODO: #1]}}}
\newcommand{\result}[1]{\textcolor{blue}{\textbf{[RESULT: #1]}}}
\newcommand{\bpb}{\mathrm{BPB}}
\newcommand{\deficit}{\mathcal{D}}
\newcommand{\gain}{\mathcal{G}}
\newcommand{\closure}{\mathcal{C}}

\title{%
  When Does Memory Pay Off?\\
  Diagnosing Positional Context Deficit in Restricted-Attention Transformers%
}
\author{Robby Sneiderman}
\date{Draft --- \today}

\begin{document}
\maketitle

% =====================================================================
\begin{abstract}
% =====================================================================
We introduce a diagnostic framework for evaluating when explicit memory
mechanisms justify their overhead in context-restricted transformers.
Standard evaluation reports a single aggregate perplexity, which averages
over positions where local attention is already sufficient and positions
where it is not. We decompose performance by position, defining a
\emph{positional context deficit} as the per-position gap between a
local-attention model and a matched global-attention reference. We then
ask whether persistent memory preferentially closes this deficit as
sequence length grows. In controlled experiments on a 12-layer GPT
trained at context lengths 2048, 4096, and 8192, we find that:
\result{deficit grows / is flat with context},
\result{memory closes X\% of deficit in high-deficit region},
and \result{aggregate BPB crosses over at T $\approx$ X / never crosses}.
The framework is architecture-agnostic and provides a reusable protocol
for diagnosing whether any memory mechanism is solving the right problem.
Code and results are available at
\url{https://github.com/Robby955/memory-bench}.
\end{abstract}


% =====================================================================
\section{Introduction}
\label{sec:intro}
% =====================================================================

Efficient attention patterns---sliding windows, sparse attention,
mixed local-global stacks---reduce the quadratic cost of full attention
at the price of restricting each token's direct access to context
\citep{jiang2023mistral, beltagy2020longformer, zaheer2020bigbird}.
Memory mechanisms such as persistent memory tokens
\citep{sukhbaatar2019persistent}, recurrent memory
\citep{bulatov2022rmt}, and test-time training
\citep{sun2024ttt} are motivated as a way to compensate: they
maintain state across positions or segments, potentially recovering
information that restricted attention cannot reach directly.

But evaluation of these mechanisms typically reports a single aggregate
metric---perplexity or bits-per-byte (BPB)---over all positions in the
sequence. Aggregate BPB alone cannot distinguish between three cases:
(i)~there is no meaningful access problem,
(ii)~there is an access problem but memory does not repair it, and
(iii)~there is an access problem and memory repairs it enough to matter.
\citet{liu2023lost} showed that model performance depends heavily on
where relevant information sits in the context.
\citet{fang2025longppl} argued that ordinary perplexity dilutes
long-context ability by averaging over tokens that are not
context-sensitive. This study is designed to separate those cases.

We propose a position-resolved evaluation framework built on three
metrics:

\begin{enumerate}
\item \textbf{Positional context deficit} $\deficit_T(p)$: the
  per-position BPB gap between a local-attention model and a matched
  global-attention reference. This measures the cost of restricting
  context access at each position.

\item \textbf{Memory gain} $\gain_T(p)$: the per-position BPB
  improvement from adding a memory mechanism to the local-attention
  model.

\item \textbf{Deficit closure} $\closure_T(p)$: the fraction of the
  context deficit recovered by memory. This answers: where the
  architecture is impaired, does memory repair the impairment?
\end{enumerate}

The framework separates two questions that aggregate BPB conflates:
(a)~does memory help \emph{somewhere}, and (b)~does it help
\emph{where the base architecture is most deficient}? A mechanism
that provides uniform overhead reduction is less interesting than one
that preferentially targets the access deficit.

We apply this framework to persistent memory
\citep{sukhbaatar2019persistent} in a 12-layer GPT with mixed
local-global attention (SSSL pattern), training at context lengths
2048, 4096, and 8192 on a fixed token budget. This is a case study;
the framework itself is architecture-agnostic.

\paragraph{Contributions.}
\begin{itemize}
\item A reusable diagnostic protocol---deficit, gain, closure---for
  evaluating memory mechanisms under restricted-context attention.
\item Controlled experiments at three context lengths with a matched
  global-attention reference, including block bootstrap confidence
  intervals and a difference-in-differences analysis at the attention
  window boundary.
\item Targeted synthetic probes that test cross-window information
  transfer independent of natural-text confounds.
\item An honest assessment: persistent memory
  \result{does / does not} close the deficit, and aggregate
  BPB \result{does / does not} cross over.
\end{itemize}


% =====================================================================
\section{Framework}
\label{sec:framework}
% =====================================================================

\subsection{Setup}

We consider a family of autoregressive language models that share
architecture (depth, width, optimizer, token budget) and differ only in
attention pattern and memory mechanism. Let $\bpb_A(p, T)$ denote the
mean bits-per-byte at position bucket $p$ for model $A$ trained with
maximum context length $T$.

Three model variants define the framework:
\begin{itemize}
\item \textbf{Local}: window pattern SSSL (3/4 of layers use a short
  window of $\approx T/4$, 1/4 use full context). This is the standard
  efficient-attention baseline.
\item \textbf{Global}: window pattern L (all layers use full context).
  This is an operational reference---the same model with unrestricted
  access---not an information-theoretic ceiling.
\item \textbf{Memory}: window pattern SSSL plus a memory mechanism
  (persistent memory in our case study).
\end{itemize}

\subsection{Metrics}

\paragraph{Positional context deficit.}
\begin{equation}
  \deficit_T(p) = \bpb_{\text{local}}(p, T) - \bpb_{\text{global}}(p, T)
  \label{eq:deficit}
\end{equation}
Positive deficit means local attention underperforms relative to
unrestricted access at position~$p$. The deficit is an operational proxy,
not a claim about information sufficiency: global attention is a different
inductive bias, not a pure upper bound.

\paragraph{Memory gain.}
\begin{equation}
  \gain_T(p) = \bpb_{\text{local}}(p, T) - \bpb_{\text{memory}}(p, T)
  \label{eq:gain}
\end{equation}
Positive gain means the memory mechanism improves on the local-attention
model at position~$p$.

\paragraph{Deficit closure.}
\begin{equation}
  \closure_T(p) = \frac{\gain_T(p)}{\max\bigl(\deficit_T(p),\; \epsilon\bigr)}
  \label{eq:closure}
\end{equation}
with $\epsilon = 0.001$~BPB. Closure near 1.0 means memory fully
recovers the access deficit; near 0.0 means it does not help where
attention is impaired; greater than 1.0 means it exceeds the
global-attention reference (overshooting). Closure is undefined
(masked) where $\deficit_T(p) < \epsilon$, because the local model
is already sufficient there.

\subsection{Primary Endpoint and Selectivity}
\label{sec:primary}

The primary endpoint is \textbf{seed-averaged late-region closure
at $T{=}8192$}, where the late region $R_{\text{late}}$ is the
final quartile of sequence positions. This is pre-registered: the
region is fixed by position, not selected post hoc from the data.

\textbf{Regional closure} over a fixed region $R$:
\begin{equation}
  \closure(T, R)
  = \frac{\bar{B}_{\text{local}}(T, R) - \bar{B}_{\text{memory}}(T, R)}
         {\bar{B}_{\text{local}}(T, R) - \bar{B}_{\text{global}}(T, R)}
  \label{eq:regional_closure}
\end{equation}
where $\bar{B}_m(T, R)$ is the token-weighted average BPB for
model~$m$ in region~$R$. Closure is only interpreted when the
denominator is positive (i.e., the local model is worse than the
global reference in that region).

\textbf{Selectivity} measures whether memory is targeted or
uniformly shifts BPB:
\begin{equation}
  S(T) = \closure(T, R_{\text{late}}) - \closure(T, R_{\text{early}})
  \label{eq:selectivity}
\end{equation}
Positive selectivity means memory helps more in late positions
(where the deficit is expected to concentrate) than in early
positions.

\subsection{Hypotheses}
\label{sec:hypotheses}

\begin{description}
\item[H1 (late-position deficit):]
  The positional context deficit is larger in late positions than
  early positions at each context length.

\item[H2 (deficit grows with context):]
  The late-position deficit grows with context length:
  $\bar{\deficit}(R_{\text{late}}, 8192)
  > \bar{\deficit}(R_{\text{late}}, 4096)
  > \bar{\deficit}(R_{\text{late}}, 2048)$.

\item[H3 (positive closure):]
  Persistent memory closes a positive fraction of the late-position
  deficit: $\closure(8192, R_{\text{late}}) > 0$.

\item[H4 (selective closure):]
  Closure is selective: $S(T) > 0$ for $T \geq 4096$. Memory helps
  more in the deficit-heavy late region than in the early region.

\item[H5 (synthetic probe confirmation):]
  On targeted probes, the memory advantage appears primarily when
  the dependency distance exceeds the SSSL window ($T/4$).
\end{description}

\paragraph{Interpretation of outcomes.}
If H1 and H2 fail, local attention is already sufficient in this
regime. That is a useful negative result: it means the architecture
is not context-starved.
If H1 and H2 hold but H3 fails, the deficit is real but persistent
memory is not the right repair.
If H3 holds but aggregate BPB still loses, persistent memory solves
the right problem but too expensively.
If H3 holds and aggregate BPB improves, that is the strongest outcome.


% =====================================================================
\section{Experimental Design}
\label{sec:design}
% =====================================================================

\subsection{Architecture}

All models use a 12-layer GPT (nanochat codebase) with rotary position
embeddings (RoPE, $\theta = 100{,}000$), trained with Muon~+~AdamW
on FineWeb-Edu. The total token budget is fixed at 524,288 tokens per
batch across all context lengths.

\paragraph{Attention patterns.}
The SSSL pattern assigns 3/4 of layers a short window of
$\lceil T/4 / 128 \rceil \times 128$ tokens (aligned to FlashAttention-3
tile size) and 1/4 of layers the full context. Pattern~L gives all layers
the full context. Both patterns are set via a single configuration flag;
no architectural changes are needed.

\paragraph{Persistent memory.}
Following \citet{sukhbaatar2019persistent}, we prepend $K$ learnable
memory tokens to the key-value cache of designated attention layers.
These tokens attend to and are attended by all positions, providing a
fixed-capacity persistent store accessible from any position in the
sequence. The mechanism adds $K \times d_{\text{model}}$ parameters
per layer.

\subsection{Run Matrix}

\begin{table}[h]
\centering
\caption{Experiment matrix. Each cell is one training run. We reuse
existing 2048 results (4 seeds) and run new experiments at 4096 and
8192 (3 seeds each). Total new runs: 21.}
\label{tab:matrix}
\begin{tabular}{llccl}
\toprule
Context & Model & Window & Seeds & Status \\
\midrule
2048 & Local baseline & SSSL & 4 & Done (v1) \\
2048 & Global reference & L & 3 & New \\
2048 & Persistent memory & SSSL & 4 & Done (v1) \\
\midrule
4096 & Local baseline & SSSL & 3 & New \\
4096 & Global reference & L & 3 & New \\
4096 & Persistent memory & SSSL & 3 & New \\
\midrule
8192 & Local baseline & SSSL & 3 & New \\
8192 & Global reference & L & 3 & New \\
8192 & Persistent memory & SSSL & 3 & New \\
\bottomrule
\end{tabular}
\end{table}

\subsection{Controls and Admitted Confounds}

Architecture, optimizer, hyperparameters, and total token budget are
held constant. Context length is the primary independent variable;
window pattern (SSSL vs L) and mechanism (none vs persistent) are
secondary.

\textbf{Admitted confounds.} Varying context length also changes:
number of optimization steps (same token budget, larger sequences =
fewer steps), gradient noise profile, data packing statistics, and
the absolute SSSL window size ($T/4$ at $T{=}8192$ is 2048 tokens
vs 512 at $T{=}2048$). These are documented as limitations, not
controlled away.


% =====================================================================
\section{Statistical Analysis}
\label{sec:statistics}
% =====================================================================

Adjacent position buckets share overlapping context and underlying
text continuity. Naive significance tests on bucketed BPB overstate
significance. We handle autocorrelation with three complementary
methods and report all three.

\paragraph{Block bootstrap (primary).}
We resample contiguous blocks of position buckets, preserving
local autocorrelation structure. Block size is chosen as the
smallest lag~$k$ where the autocorrelation function of the deficit
series drops below 0.05, with a floor of 3 and cap at $N/4$
\citep[following][]{metaasr}. We report 95\% percentile
confidence intervals from 5,000 bootstrap samples for mean
deficit, gain, and closure.

\paragraph{Newey-West HAC standard errors.}
When testing whether deficit varies with position via OLS, we use
Newey-West heteroskedasticity and autocorrelation consistent (HAC)
standard errors \citep{newey1987hac} with a Bartlett kernel and
automatic lag selection ($\lfloor N^{1/3} \rfloor$).

\paragraph{Effective degrees of freedom.}
As a simple check, we compute effective sample size using the
\citet{bretherton1999effective} formula:
$N_{\text{eff}} = N \cdot (1 - r_1) / (1 + r_1)$, where $r_1$
is the lag-1 autocorrelation.

\paragraph{Difference-in-differences at the window boundary.}
We partition position buckets into Near ($p < T/4$, within the
SSSL short window for all layers) and Far ($p \geq T/4$, beyond
the short-window boundary). The DID estimator is:
\begin{equation}
  \text{DID} = \bigl[\bpb_{\text{mem}}(\text{far}) - \bpb_{\text{local}}(\text{far})\bigr]
             - \bigl[\bpb_{\text{mem}}(\text{near}) - \bpb_{\text{local}}(\text{near})\bigr]
  \label{eq:did}
\end{equation}
If memory specifically helps beyond the window boundary,
$\text{DID} < 0$. The parallel trends assumption is testable:
we check that local and memory models have similar
position-dependent BPB trends in the Near region.


% =====================================================================
\section{Results}
\label{sec:results}
% =====================================================================

\subsection{Deficit Maps}

\result{Figure~\ref{fig:deficit_map}: three panels showing
local vs global BPB by position at each context length, with
shaded deficit region and T/4 boundary marked. Expected: deficit
grows in late positions as T increases.}

\begin{figure}[h]
\centering
% \includegraphics[width=\textwidth]{figures/fig_deficit_map.png}
\fbox{\parbox{0.95\textwidth}{\centering\vspace{2cm}
  \textbf{Figure placeholder: Deficit maps at T=2048, 4096, 8192}\\
  Local (SSSL) vs Global (L) BPB by position, shaded deficit region.
\vspace{2cm}}}
\caption{Positional context deficit at three context lengths.
  \textit{Blue}: local-attention baseline (SSSL).
  \textit{Green}: global-attention reference (L).
  \textit{Red shading}: context deficit (local $-$ global).
  Dashed line: SSSL window boundary at $T/4$.
  \result{Describe what we observe.}}
\label{fig:deficit_map}
\end{figure}

\subsection{Deficit Closure}

\result{Figure~\ref{fig:closure}: deficit and gain curves with
closure in bottom panels. Table~\ref{tab:closure}: primary
endpoint values.}

\begin{table}[h]
\centering
\caption{Primary endpoint: mean deficit closure in the late region
(final quartile of sequence positions), with 95\% block bootstrap
confidence intervals.}
\label{tab:closure}
\begin{tabular}{lccc}
\toprule
Context & Late Deficit & $\closure(T, R_{\text{late}})$ & 95\% CI \\
\midrule
2048 & \result{X.XXXXX} & \result{X.XXX} & \result{[X.XX, X.XX]} \\
4096 & \result{X.XXXXX} & \result{X.XXX} & \result{[X.XX, X.XX]} \\
8192 & \result{X.XXXXX} & \result{X.XXX} & \result{[X.XX, X.XX]} \\
\bottomrule
\end{tabular}
\end{table}

\begin{figure}[h]
\centering
\fbox{\parbox{0.95\textwidth}{\centering\vspace{2cm}
  \textbf{Figure placeholder: Deficit closure panels}\\
  Top row: deficit (red) and gain (green) curves.
  Bottom row: closure curve with mean annotation.
\vspace{2cm}}}
\caption{Deficit closure for persistent memory at three context lengths.
  \textit{Top}: per-position deficit and gain.
  \textit{Bottom}: closure (gain / deficit). Dashed line at
  $\closure = 1.0$ (full closure).
  \result{Describe trend.}}
\label{fig:closure}
\end{figure}

\subsection{Difference-in-Differences}

\begin{table}[h]
\centering
\caption{Difference-in-differences at the SSSL window boundary ($T/4$).
Negative DID indicates memory helps more beyond the boundary.}
\label{tab:did}
\begin{tabular}{lcccc}
\toprule
Context & Near Effect & Far Effect & DID & Parallel Trends \\
\midrule
2048 & \result{$\pm$X.XXXXX} & \result{$\pm$X.XXXXX} & \result{$\pm$X.XXXXX} & \result{yes/no} \\
4096 & \result{$\pm$X.XXXXX} & \result{$\pm$X.XXXXX} & \result{$\pm$X.XXXXX} & \result{yes/no} \\
8192 & \result{$\pm$X.XXXXX} & \result{$\pm$X.XXXXX} & \result{$\pm$X.XXXXX} & \result{yes/no} \\
\bottomrule
\end{tabular}
\end{table}

\subsection{Synthetic Probes}

We verify that the natural-text deficit arises from cross-window
information access by testing three controlled probes:

\begin{enumerate}
\item \textbf{Token recall at distance}: A novel numeric association
  (``Item alpha-\{id\} costs \{value\}'') is planted at the sequence
  start, followed by natural-text filler, then queried at varying
  distances. We measure prediction loss on the value tokens, sweeping
  distance across the $T/4$ window boundary.

\item \textbf{Entity tracking}: Same structure with entity-attribute
  bindings (``\{Name\} works at \{Org\}''), using fictional entities
  to prevent training-data memorization.

\item \textbf{Cross-boundary associative recall}: Multiple key-value
  pairs are planted in the first $T/4$ of the sequence and queried
  in the second half, testing multi-item recall capacity across the
  window boundary.
\end{enumerate}

\result{Figure or table showing probe loss by distance.
Expected: memory models show lower loss at distances $> T/4$,
confirming cross-window information transfer.}

\subsection{Aggregate BPB}

For completeness, we report overall BPB:

\begin{table}[h]
\centering
\caption{Aggregate BPB (mean $\pm$ std across seeds).
  Persistent memory \result{does/does not} achieve lower aggregate
  BPB than baseline at context $\geq$ \result{X}.}
\label{tab:aggregate}
\begin{tabular}{lccc}
\toprule
Context & Local (SSSL) & Global (L) & Persistent Memory \\
\midrule
2048 & $0.8459 \pm 0.0006$ & \result{X.XXXX $\pm$ X.XXXX} & $0.8442 \pm 0.0005$ \\
4096 & \result{X.XXXX} & \result{X.XXXX} & \result{X.XXXX} \\
8192 & \result{X.XXXX} & \result{X.XXXX} & \result{X.XXXX} \\
\bottomrule
\end{tabular}
\end{table}

\paragraph{Interpretation.}
The aggregate metric is reported for calibration but is not the
primary evidence. As \citet{fang2025longppl} argue, aggregate
perplexity can hide long-context ability by averaging over tokens
that are not context-sensitive. The deficit and closure metrics
provide a sharper diagnostic.


% =====================================================================
\section{Related Work}
\label{sec:related}
% =====================================================================

\paragraph{Position-dependent evaluation.}
\citet{liu2023lost} established that LLM performance varies with the
position of relevant information, finding a U-shaped accuracy curve
across context. \citet{kamradt2024niah} introduced needle-in-a-haystack
position $\times$ context heatmaps using synthetic probes.
\citet{liu2024forgetting} measured position-resolved memorization
accuracy, showing qualitatively different forgetting curves for
Transformers and SSMs. \citet{fang2025longppl} decomposed perplexity by
token type (context-sensitive vs not), arguing that standard perplexity
hides long-context ability. Our work decomposes by \emph{position} and
introduces the deficit/closure framework, which is complementary to
LongPPL's token-type decomposition.

\paragraph{Theoretical grounding.}
\citet{xiao2025sliding} showed that effective influence in stacked
sliding windows decays exponentially with the number of hops, not
linearly, making the theoretical receptive field vastly larger than the
empirical one. \citet{chowdhury2026lost} demonstrated that the U-shaped
position bias exists at initialization before training, arising from
causal masking and residual connections. \citet{wu2025position} and
\citet{herasimchyk2026residual} provide formal proofs of U-shaped
attention concentration at finite depth due to residual connection
barriers. These results provide the theoretical mechanism behind the
positional deficit we measure.

\paragraph{Memory mechanisms.}
Persistent memory \citep{sukhbaatar2019persistent} prepends learnable
tokens to the key-value cache, providing cross-position information
transfer at fixed parameter cost. Recurrent Memory Transformer
\citep{bulatov2022rmt} processes segments sequentially with read/write
memory operations. Test-time training \citep{sun2024ttt} updates model
parameters at inference. Gated DeltaNet \citep{yang2025gated} replaces
attention with a linear recurrence. These mechanisms are all motivated
as compensating for limited context access; our framework provides a
controlled way to test that motivation.

\paragraph{Restricted-context attention.}
Mistral~7B \citep{jiang2023mistral} uses a sliding window of 4096
tokens. Longformer \citep{beltagy2020longformer} combines local
windowed attention with task-specific global tokens.
BigBird \citep{zaheer2020bigbird} uses a mix of random, window,
and global attention. Our SSSL pattern (from the nanochat codebase)
assigns most layers a short window and a minority the full context,
similar in spirit to these production architectures.


% =====================================================================
\section{Discussion}
\label{sec:discussion}
% =====================================================================

\paragraph{What we showed.}
\result{1--2 sentences summarizing the main finding.
Expected: restricted attention creates a measurable positional
deficit that grows with context, and persistent memory
closes / does not close a meaningful fraction of it.}

\paragraph{What we did not show.}
We did not prove that memory is necessary for long context in general.
The global-attention reference is an operational proxy, not an
information-theoretic ceiling. Correlation between deficit and gain
does not establish causation; the synthetic probes provide controlled
evidence but not proof. The results are specific to one architecture
family, one training regime, and one memory mechanism.

\paragraph{Limitations.}
\begin{itemize}
\item \textbf{Scale.} All models are 12-layer ($\sim$50M parameters).
  Results may not transfer to larger scales where attention patterns
  and information routing differ qualitatively.
\item \textbf{Confounds.} Varying context length changes step count,
  gradient noise, and data packing. These are documented but not
  controlled.
\item \textbf{Single mechanism.} We study persistent memory only.
  Other mechanisms (RMT, TTT, DeltaNet) may show different
  closure profiles.
\item \textbf{Token averaging.} Our bucketed BPB averages over all
  tokens in each position bucket, including tokens that may not be
  context-sensitive \citep{fang2025longppl}. A key-token analysis
  would sharpen the signal.
\end{itemize}

\paragraph{Broader applicability.}
The deficit/closure framework is not specific to persistent memory or
the SSSL attention pattern. Any architecture with restricted context
access can be evaluated against a matched unrestricted reference.
Builders of sparse-attention models, hierarchical memory systems, or
retrieval-augmented context compressors can reuse this protocol to
diagnose whether their mechanism is solving an access deficit or
adding overhead to positions that are already sufficient.


% =====================================================================
\section{Conclusion}
\label{sec:conclusion}
% =====================================================================

We introduced positional context deficit maps as a diagnostic for
context-restricted transformers and applied them to evaluate persistent
memory at three context lengths. The framework separates two questions
that aggregate BPB conflates: whether memory helps somewhere, and
whether it helps where the base architecture is most deficient.
\result{One sentence with the headline result.}
The protocol is reusable, reproducible, and informative even under
null results: a small deficit means the architecture is not
context-starved; zero closure means memory is not the right
intervention; high closure means memory is paying for itself.
Code and all results are available at
\url{https://github.com/Robby955/memory-bench}.


% =====================================================================
% References
% =====================================================================
\bibliographystyle{plainnat}
\bibliography{references}

\end{document}
